% GENERATED FILE. Edit canonical lessons, then run npm run generate:books.
% canonical-source-hash: fnv1a64:83965d414a5711da
% canonical-lessons: ZH-C36-lai, ZH-W36-lai, ZH-C36-meiguo, ZH-C36-yingguo, ZH-C36-beijing, ZH-C36-na, ZH-R36-review-31-36

\chapter{Where People Come From}
\label{ch:zh-where-people-come-from}
\hlchaptermodality{36}

\begin{chapteropening}
\textbf{By the end of this chapter:} \emph{I can say where somebody comes from and ask which country.}

Complete the last lesson of chapter 36: I can say where somebody comes from and ask which country.
\end{chapteropening}

\section[\texorpdfstring{\zh{来} (lái)}{来 (lái)}]{\zh{来} (lái) — to come}

\label{lesson:ZH-C36-lai}



\begin{quote}
Before the new one: say the Mandarin for sir, then the Mandarin for to live somewhere.
\end{quote}

\subsection*{You'll want to know: \zh{来}}
\textbf{\zh{来}} — \emph{lái} — \textquotedblleft{}to come\textquotedblright{}.

\begin{quote}
\textbf{\zh{妈妈来了}} \emph{māma lái le} — Mum's here.
\end{quote}

The new character, in pieces: \zh{来} — one piece built from the top down: a short bar, a dot and a short stroke under it, a long bar, a vertical through it, and two strokes spreading at the foot. The next lesson puts it in your hand.

\begin{cousinweb}
\textbf{\zh{来}} (\emph{lái}) is inherited from Middle Chinese \emph{loj} (Guangyun, in Baxter's transcription), glossed 'to arrive'.
\end{cousinweb}

\begin{grammarlens}[title={What you've built}]
One more word about people.
\end{grammarlens}

\subsection*{Your turn}
\begin{itemize}
\raggedright
  \item \emph{Say it:} \emph{lái}
  \item \emph{Say it:} \emph{lái}, once more
  \item \emph{Say it:} \emph{māma lái le} — Mum's here.
  \item \emph{Recall:} say the Mandarin for sir, then the Mandarin for to live somewhere, then say \emph{lái} again
\end{itemize}

\begin{tcolorbox}[breakable,colback=teal!4,colframe=teal!35!black,title={Before you move on}]
What does \zh{来} mean? (\textquotedblleft{}To come\textquotedblright{}.) Say it once more.
\end{tcolorbox}
\section[\texorpdfstring{\zh{来} (lái)}{来 (lái)}]{\zh{来} — one piece, top to bottom}

\label{lesson:ZH-W36-lai}



\begin{quote}
The word first, then the character inside it.

\begin{itemize}
\raggedright
  \item \emph{Recall:} say \emph{lái} — \textbf{R1}, one lesson back, before it decays
\end{itemize}

\begin{quote}
\zh{大} \zh{十}
\end{quote}

\emph{Read it:} each one aloud, with nothing said first — \emph{dà}, \emph{shí}
\end{quote}

\subsection*{Script — one piece, top to bottom}
\hlblockfigure{\detokenize{figures/ZH-W36-lai-filmstrip.pdf}}{How \zh{来} is written, stroke by stroke}

\begin{quote}
\zh{来}
\end{quote}

\textbf{One piece, three layers, top to bottom.}

\textbf{The top layer, three strokes.} A short horizontal. Under it, a dot on the left and a short left-falling stroke on the right, leaning in toward each other.

\textbf{The middle, two strokes.} A long horizontal, then the vertical through the middle of everything, from the very top, through both bars, to the foot.

\textbf{The foot, two strokes.} A left-falling stroke and a right-falling stroke, spreading out from the vertical like the legs of \zh{大}.

3 + 2 + 2 = seven.

\subsection*{Your turn}
\begin{itemize}
\raggedright
  \item \emph{Write it:} \zh{来}, all seven, with the model in front of you
  \item \emph{Cover:} write it again from memory, and count the pieces, not the strokes
\end{itemize}

\begin{tcolorbox}[breakable,colback=teal!4,colframe=teal!35!black,title={Before you move on}]
How many layers in \zh{来}? (\textbf{Three}.) Which comes first, the dot or the short left-falling stroke? (\textbf{The dot}.) Which character you know ends with the same two spreading strokes? (\textbf{\zh{大}}.)

Source: \href{https://raw.githubusercontent.com/chanind/hanzi-writer-data/68d10a4b21150cae5e1ebbd223eed289cf32d90c/data/\%E6\%9D\%A5.json}{Hanzi Writer Data \zh{来}, pinned stroke paths and medians}.
\end{tcolorbox}
\section[Měiguó]{Měiguó — America; the United States}

\label{lesson:ZH-C36-meiguo}



\begin{quote}
Before the new one: say the Mandarin for to live somewhere, then the Mandarin for to come.
\end{quote}

\subsection*{You'll want to know: Měiguó}
\emph{Měiguó} — \textquotedblleft{}America; the United States\textquotedblright{}.

\begin{quote}
\emph{Měiguó hěn dà} — America is huge.
\end{quote}

Its characters come later; for now it is a word for the ear.

\begin{cousinweb}
\emph{Měiguó} is built from měi 'beautiful' and \zh{国} 'country', both inherited from Middle Chinese, \emph{mijX} and \emph{kwok} (Guangyun, in Baxter's transcription). Měi is usually explained as standing for a syllable of the name America.
\end{cousinweb}

\begin{grammarlens}[title={What you've built}]
One more word about people.
\end{grammarlens}

\subsection*{Your turn}
\begin{itemize}
\raggedright
  \item \emph{Say it:} \emph{Měiguó}
  \item \emph{Say it:} \emph{Měiguó}, once more
  \item \emph{Say it:} \emph{Měiguó hěn dà} — America is huge.
  \item \emph{Recall:} say the Mandarin for to live somewhere, then the Mandarin for to come, then say \emph{Měiguó} again
\end{itemize}

\begin{tcolorbox}[breakable,colback=teal!4,colframe=teal!35!black,title={Before you move on}]
What does Měiguó mean? (\textquotedblleft{}America; the United States\textquotedblright{}.) Say it once more.
\end{tcolorbox}
\section[Yīngguó]{Yīngguó — Britain; the United Kingdom}

\label{lesson:ZH-C36-yingguo}



\begin{quote}
Before the new one: say the Mandarin for to come, then the Mandarin for the United States.
\end{quote}

\subsection*{You'll want to know: Yīngguó}
\emph{Yīngguó} — \textquotedblleft{}Britain; the United Kingdom\textquotedblright{}.

\begin{quote}
\emph{Yīngguó hěn xiǎo} — Britain is quite small.
\end{quote}

Its characters come later; for now it is a word for the ear.

\begin{cousinweb}
\emph{Yīngguó} is built from yīng 'flower, fine' and \zh{国} 'country', both inherited from Middle Chinese, \emph{'jaeng} and \emph{kwok} (Guangyun, in Baxter's transcription). Yīng is usually explained as standing for the first syllable of England.
\end{cousinweb}

\begin{grammarlens}[title={What you've built}]
One more word about people.
\end{grammarlens}

\subsection*{Your turn}
\begin{itemize}
\raggedright
  \item \emph{Say it:} \emph{Yīngguó}
  \item \emph{Say it:} \emph{Yīngguó}, once more
  \item \emph{Say it:} \emph{Yīngguó hěn xiǎo} — Britain is quite small.
  \item \emph{Recall:} say the Mandarin for to come, then the Mandarin for the United States, then say \emph{Yīngguó} again
\end{itemize}

\begin{tcolorbox}[breakable,colback=teal!4,colframe=teal!35!black,title={Before you move on}]
What does Yīngguó mean? (\textquotedblleft{}Britain; the United Kingdom\textquotedblright{}.) Say it once more.
\end{tcolorbox}
\section[Běijīng]{Běijīng — Beijing}

\label{lesson:ZH-C36-beijing}



\begin{quote}
Before the new one: say the Mandarin for the United States, then the Mandarin for the United Kingdom.
\end{quote}

\subsection*{You'll want to know: Běijīng}
\emph{Běijīng} — \textquotedblleft{}Beijing\textquotedblright{}.

\begin{quote}
\emph{wǒmen lái Běijīng le} — We've come to Beijing.
\end{quote}

Its characters come later; for now it is a word for the ear.

\begin{cousinweb}
\emph{Běijīng} is built from běi 'north' and jīng 'capital', both inherited from Middle Chinese, \emph{pok} and \emph{kjaeng} (Guangyun, in Baxter's transcription): the northern capital.
\end{cousinweb}

\begin{grammarlens}[title={What you've built}]
One more word about people.
\end{grammarlens}

\subsection*{Your turn}
\begin{itemize}
\raggedright
  \item \emph{Say it:} \emph{Běijīng}
  \item \emph{Say it:} \emph{Běijīng}, once more
  \item \emph{Say it:} \emph{wǒmen lái Běijīng le} — We've come to Beijing.
  \item \emph{Recall:} say the Mandarin for the United States, then the Mandarin for the United Kingdom, then say \emph{Běijīng} again
\end{itemize}

\begin{tcolorbox}[breakable,colback=teal!4,colframe=teal!35!black,title={Before you move on}]
What does Běijīng mean? (\textquotedblleft{}Beijing\textquotedblright{}.) Say it once more.
\end{tcolorbox}
\section[nǎ]{nǎ — which}

\label{lesson:ZH-C36-na}



\begin{quote}
Before the new one: say the Mandarin for the United Kingdom, then the Mandarin for Beijing.
\end{quote}

\subsection*{You'll want to know: nǎ}
\emph{nǎ} — \textquotedblleft{}which\textquotedblright{}.

\begin{quote}
\emph{nǐ shì nǎ guó rén} — Which country are you from?
\end{quote}

Its characters come later; for now it is a word for the ear.

\begin{cousinweb}
\emph{nǎ} has no older form in any dictionary we use, so its history is not settled.
\end{cousinweb}

\begin{grammarlens}[title={What you've built}]
One more word about people.
\end{grammarlens}

\subsection*{Your turn}
\begin{itemize}
\raggedright
  \item \emph{Say it:} \emph{nǎ}
  \item \emph{Say it:} \emph{nǎ}, once more
  \item \emph{Say it:} \emph{nǐ shì nǎ guó rén} — Which country are you from?
  \item \emph{Recall:} say the Mandarin for the United Kingdom, then the Mandarin for Beijing, then say \emph{nǎ} again
\end{itemize}

\begin{tcolorbox}[breakable,colback=teal!4,colframe=teal!35!black,title={Before you move on}]
What does nǎ mean? (\textquotedblleft{}Which\textquotedblright{}.) Say it once more.
\end{tcolorbox}
\section[(dialogue)]{Looking back: from He, She and We to Where People Come From}

\label{lesson:ZH-R36-review-31-36}



\begin{quote}
Say the Mandarin for she and he.
\end{quote}

\subsection*{Your turn}
\begin{itemize}
\raggedright
  \item \emph{Say it:} she, he, they, we, you all
  \item \emph{Say it:} mum, dad, elder brother, elder sister, younger brother
  \item \emph{Say it:} male, man, woman, younger sister, child
  \item \emph{Say it:} to be called, you (polite), to have the family name, to know someone, friend
  \item \emph{Say it:} to work, old, teacher, sir, to live somewhere
  \item \emph{Say it:} to come, the United States, the United Kingdom, Beijing, which
  \item \emph{Read it:} \zh{他} \zh{们} \zh{妈} \zh{男} \zh{叫} \zh{来} — each one aloud, then a word that holds it
\end{itemize}

\begin{tcolorbox}[breakable,colback=teal!4,colframe=teal!35!black,title={Before you move on}]
She? (\emph{tā}.) He? (\textbf{\zh{他}}, \emph{tā}.) They? (\textbf{\zh{他们}}, \emph{tāmen}.) We? (\textbf{\zh{我们}}, \emph{wǒmen}.) You all? (\textbf{\zh{你们}}, \emph{nǐmen}.) Mum? (\textbf{\zh{妈妈}}, \emph{māma}.) Dad? (\emph{bàba}.) Elder brother? (\emph{gēge}.) Elder sister? (\emph{jiějie}.) Younger brother? (\emph{dìdi}.) Male? (\textbf{\zh{男}}, \emph{nán}.) Man? (\textbf{\zh{男人}}, \emph{nánrén}.) Woman? (\textbf{\zh{女人}}, \emph{nǚrén}.) Younger sister? (\emph{mèimei}.) Child? (\emph{háizi}.) To be called? (\textbf{\zh{叫}}, \emph{jiào}.) You (polite)? (\emph{nín}.) To have the family name? (\emph{xìng}.) To know someone? (\emph{rènshi}.) Friend? (\emph{péngyou}.) To work? (\emph{gōngzuò}.) Old? (\emph{lǎo}.) Teacher? (\emph{lǎoshī}.) Sir? (\emph{xiānsheng}.) To live somewhere? (\emph{zhù}.) To come? (\textbf{\zh{来}}, \emph{lái}.) The United States? (\emph{Měiguó}.) The United Kingdom? (\emph{Yīngguó}.) Beijing? (\emph{Běijīng}.) Which? (\emph{nǎ}.)
\end{tcolorbox}
